\documentclass[11pt]{article}
\usepackage[margin=1in]{geometry}
\usepackage{amsmath,amssymb,amsthm,mathtools}
\usepackage[T1]{fontenc}
\usepackage{lmodern,microtype}
\usepackage{xurl}
\usepackage[colorlinks=true,linkcolor=blue,citecolor=blue,urlcolor=blue]{hyperref}
\hypersetup{pdftitle={A conditional obstruction to polynomial sample complexity in unrestricted NNbayes},pdfauthor={ChatGPT-6 Astra; Kaarel Haenni; Dmitry Vaintrob}}
\newtheorem{theorem}{Theorem}[section]
\newtheorem{lemma}[theorem]{Lemma}
\newtheorem{proposition}[theorem]{Proposition}
\newtheorem{corollary}[theorem]{Corollary}
\newtheorem{definition}[theorem]{Definition}
\theoremstyle{remark}\newtheorem{remark}[theorem]{Remark}
\newcommand{\R}{\mathbb R}
\newcommand{\E}{\mathbb E}
\newcommand{\Pp}{\mathsf P/\mathrm{poly}}
\newcommand{\TC}{\mathsf{TC}}
\newcommand{\NC}{\mathsf{NC}}
\newcommand{\Prb}{\mathbb P}
\newcommand{\poly}{\operatorname{poly}}
\newcommand{\sgn}{\operatorname{sign}}
\title{A conditional obstruction to polynomial sample complexity\newline in unrestricted NNbayes}
\author{ChatGPT-6 Astra \and Kaarel H\"anni \and Dmitry Vaintrob\thanks{Authors are listed in alphabetical order.}}
\date{Research draft --- 26 October 2026}
\begin{document}
\maketitle

\begin{abstract}
Assume that polynomial-size Boolean circuits cannot all be replaced by
polynomial-size, logarithmic-depth threshold circuits:
$\Pp\ne\TC^1$, with both classes nonuniform.
Then the original Gaussian tanh Bayesian learner has no distribution-free
sample bound jointly polynomial in input dimension, teacher width, and
teacher depth that is independent of learner width. This holds separately
for every prescribed mean-field-or-stronger scale schedule, including each
fixed exponent $\gamma\ge1$. The failure already occurs at fixed accuracy
and confidence. Training is noiseless; the posterior fits every observed
label exactly and is conditioned to have an absolute output margin at
every Boolean input. No computational limitation is imposed on inference.

The proof combines an established posterior circuit ceiling with a finite
majority argument: distribution-free learning by mixtures of small
threshold circuits forces an exact small threshold circuit for the target.
We give this statistical reduction and the bounded-weight Boolean-to-tanh
implementation explicitly. The analytic posterior ceiling is stated with
its full quantifiers and cited as an input. The conclusion rules out a
universal polynomial guarantee; it does not establish exponential depth
dependence. Obtaining that guarantee while retaining the target class and
the other learning conventions therefore requires changing the prior;
restricting its support is one possible change. A Lean formalization checks
the concrete-model reduction conditional on the analytic ceiling; the proof
of that analytic input remains outside the formalization.
\end{abstract}

\section{Introduction}

A natural goal for Bayesian neural-network learning is a sample bound
polynomial in the size of a circuit implementing the target, independently
of the much larger network used by the learner. At each fixed depth,
Gaussian tanh priors at mean-field scaling and stronger scales admit such
width-independent bounds. The degree or constants of those bounds can
grow rapidly with depth. It is therefore important to distinguish a
polynomial bound \emph{at every fixed depth} from one polynomial jointly
in width and depth.

This paper gives a conditional negative answer to the latter question for
the original unrestricted Gaussian prior. Under the circuit-complexity
conjecture $\Pp\ne\TC^1$, a universal joint-polynomial sample bound is
impossible. This is a statistical obstruction, even for exact Bayesian
inference with unlimited computation. It does not arise from the difficulty
of running a posterior-sampling algorithm.

The mechanism is a mismatch between target expressivity and posterior
expressivity at limited sample counts. Every small Boolean circuit can be
implemented by a polynomial-resource tanh teacher with bounded weights.
An analytic theorem for the original posterior nevertheless places most
posterior mass on functions having threshold circuits of size polynomial
in dimension, sample count, and learner depth, and of logarithmic
threshold depth. If the learner succeeded for every input distribution
using only polynomially many examples, an elementary boosting argument would
turn these posterior mixtures into an exact $\TC^1$ circuit for every
function in $\Pp$.

The width quantifier is essential: the guarantee must cover arbitrarily
large finite learner widths. So is the distribution quantifier: training
and testing use the same distribution, but the guarantee must hold for
every such distribution. No distribution shift is involved.

We present the conditional impossibility as the main result. The
underlying analytic circuit ceiling was established in
\cite{barriers}. Its proof is substantially longer than the reduction
here; Section~\ref{sec:analytic} states exactly what is imported and
explains its role. All subsequent probabilistic and complexity-theoretic
steps are supplied below.

\section{Learner, target complexity, and circuit classes}
\label{sec:model}

\subsection{The original Gaussian tanh posterior}

Let $d\ge1$, $X_d=\{-1,1\}^d$, and $f:X_d\to\{-1,1\}$.
A network has weighted depth $L\ge3$, hidden width $N$, augmented input
$h^0(x)=(1,x)$, hidden states
\[
 h^j(x)=\tanh(W^j h^{j-1}(x)),\qquad 1\le j<L,
\]
and scalar linear output $z_W(x)=W^L h^{L-1}(x)$.
The activation acts coordinatewise. There are no separate hidden biases;
constant wires use ordinary neurons and are counted in the width.

For $0<t\le1$, let $\Pi_{d,L,N,t}$ be the product Gaussian law on all
weights, with mean zero and entrywise variances
\begin{equation}\label{eq:prior}
 \frac{t^2}{d}\quad\text{in the first layer},\qquad
 \frac{t^2}{N}\quad\text{in intermediate layers},\qquad
 \frac{t^2}{N^2}\quad\text{in the readout}.
\end{equation}
The critical, or mean-field, scale is $t=1$.
The usual exponent family $\gamma\ge1$ is obtained by setting
$t=N^{(1-\gamma)/2}$.

Fix a prescribed schedule
\begin{equation}\label{eq:schedule}
 t=t(d,L,N)\in(0,1].
\end{equation}
It may depend on these architectural dimensions but not on the training
sample or the input distribution. It need not be computable.

For a labeled sample $S=((x_i,f(x_i)))_{i=1}^n$, condition the prior on
\begin{equation}\label{eq:conditioning}
 |z_W(x)|>1/2\quad(x\in X_d),\qquad
 f(x_i)z_W(x_i)>1/2\quad(1\le i\le n).
\end{equation}
Let $\Pi_S$ denote this posterior and put
\[
 p_S(y\mid x)=\Pi_S\{W:\sgn z_W(x)=y\}.
\]
For an input law $\mu$ on $X_d$, the confidence-failure mass is
\begin{equation}\label{eq:error}
 \mathcal E_\mu(f,S)
   =\mu\{x:p_S(f(x)\mid x)\le2/3\}.
\end{equation}
Thus a prediction counts as successful only when the posterior assigns
probability strictly greater than $2/3$ to the correct label.
The training inputs are iid from $\mu$, which is also the test law.

Here ``unrestricted'' means that no internal-stability or other robustness
condition is added to the prior support. The global absolute-margin
condition in \eqref{eq:conditioning} remains part of the learner.

\subsection{Teacher resources and universal polynomial learning}

For an integer $C\ge1$, a target has teacher resources $(d,C,L)$ if it admits a depth-$L$ tanh
network of hidden widths at most $C$, all weights bounded in absolute
value by $C$, and signed output greater than $1/2$ on every input.
Equivalently, its hidden widths may all be taken equal to $C$ by zero padding. Larger hidden widths can also be obtained by zero padding. Continuity on the
finite cube then gives a positive-prior neighborhood of that
implementation, so the conditional law \eqref{eq:conditioning} is defined
for every sample and every $N\ge C$.

For fixed statistical tolerances, any polynomial bound in $d,C,L$ is
bounded above by $A(1+d+C+L)^k$ for some real $A>0$ and integer $k\ge0$.
We use that envelope to state the question precisely.

\begin{definition}[Universal polynomial sample guarantee]\label{def:poly}
For a schedule \eqref{eq:schedule}, let $\mathsf{PolyLearn}(t)$ mean that
there are $A>0$ and $k\in\mathbb N$ such that, for every
$d\ge1$, $C\ge1$, $L\ge3$, every target with those teacher resources,
every $N\ge C$, every law $\mu$, and every integer sample count
\begin{equation}\label{eq:sample}
 n\ge A(1+d+C+L)^k,
\end{equation}
an iid labeled sample satisfies
\begin{equation}\label{eq:guarantee}
 \Prb_{S\sim\mu^n}\{\mathcal E_\mu(f,S)\le1/16\}\ge1/2.
\end{equation}
The constants $A,k$ are independent of learner width and input law.
\end{definition}

A polynomial guarantee at all accuracy and confidence parameters would
in particular imply \eqref{eq:guarantee}. The proof below uses the
particular integer count $n=A_0(1+d+C+L)^k$, with a fixed positive integer
$A_0\ge A$; it does not assume that posterior accuracy is monotone in $n$.

The quantity $d+C+L$ is polynomially equivalent to an ordinary upper bound
on the number of parameters in this dense teacher architecture, after
including input dimension and depth. More importantly, the reduction
below starts with ordinary Boolean circuits and uses constant-magnitude
integer weights. The obstruction does not require hiding an arbitrary
truth table in the precision of a teacher coefficient.

\subsection{Nonuniform threshold circuits}

A circuit family has one Boolean circuit for each input dimension.
``Nonuniform'' imposes no efficient procedure for generating that family.
The class $\Pp$ consists of functions with polynomial-size Boolean circuit
families. The class $\TC^1$ consists of functions with polynomial-size
threshold circuit families of depth $O(\log(d+2))$. One may use
unbounded-fan-in majority gates and ordinary Boolean gates as the
threshold basis. The class $\TC^0$ uses constant depth instead.
All circuit classes in this paper are nonuniform.

The assumption $\Pp\ne\TC^1$ says that some polynomial-size computations
cannot be parallelized into logarithmic threshold depth with polynomial
size. Appending a constant number of polynomial-fan-in majority layers
preserves $\TC^1$. The reverse simulation has no depth requirement:

\begin{lemma}[Threshold circuits are polynomial-size Boolean computations]
\label{lem:threshold-boolean}
A threshold circuit with total gate and input-wire count $s$ has an
ordinary Boolean circuit of size polynomial in $s$. Consequently
$\TC^1\subseteq\Pp$.
\end{lemma}
\begin{proof}
For a gate with $m$ incoming bits, define $T_{j,k}$ to mean that at least
$k$ of the first $j$ bits are one. The recurrence
\[
 T_{j,k}=T_{j-1,k}\ \mathsf{OR}\
             (x_j\ \mathsf{AND}\ T_{j-1,k-1})
\]
with $T_{j,0}=1$ and $T_{0,k}=0$ for $k>0$ computes every relevant count
threshold using $O(m^2)$ gates. Thresholds above $m$ are constant false.
Replacing every threshold gate this way costs at most a constant times
$\sum_j(m_j+1)^2$, which is $O((s+1)^2)$. AND, OR, and NOT have constant-size
NAND implementations. No bound on the Boolean simulation's depth is
needed for membership in $\Pp$.
\end{proof}

\section{The conditional impossibility theorem}

\begin{theorem}[No universal polynomial circuit-learning bound]
\label{thm:main}
Assume $\Pp\ne\TC^1$. For every prescribed scale schedule
$t(d,L,N)\in(0,1]$, the property $\mathsf{PolyLearn}(t)$ is false.

Equivalently, for every real $A>0$ and integer $k\ge0$, there exist dimensions
$d,C,L$, a target with those teacher resources, a learner width $N\ge C$,
a common training/test law $\mu$, and an integer sample count satisfying
\eqref{eq:sample}, such that
\begin{equation}\label{eq:failure}
 \Prb_{S\sim\mu^n}\{\mathcal E_\mu(f,S)>1/16\}>1/2.
\end{equation}
In particular the conclusion holds separately for every fixed
$\gamma\ge1$ in the exponent family.
\end{theorem}

Theorem~\ref{thm:main} concerns unlimited-computation Bayesian prediction.
Its target family may be chosen from $\Pp$, with polynomial-resource
bounded-weight teachers. The difficult distribution and the learner width
may depend on the proposed sample polynomial.

\begin{corollary}[What must change]\label{cor:change}
Under the same complexity assumption, no theorem can give the original
learner a sample guarantee polynomial jointly in $d,C,L$ at every
eligible learner width, for every target in Section~\ref{sec:model} and
every common training/test law.
Consequently, obtaining such a guarantee while retaining the architecture, that target
class, those statistical conventions, and those width and distribution
quantifiers requires changing the prior. Restricting its support to
networks satisfying an additional stability condition is one possible
change, not a conclusion that every such restriction suffices.
\end{corollary}

If other parts of the goal may change, there are other routes: restrict
the target class, assume a particular input law, limit learner width, or
alter the likelihood or prediction criterion. The theorem does not say
that a robustness restriction is the uniquely necessary remedy.

\section{The analytic input: a posterior circuit ceiling}
\label{sec:analytic}

For integers $R\ge2$, let $\mathcal H_R$ be the finite set of Boolean
functions representable by threshold circuits of size at most $R^b$ and
depth at most $b\log R$, for a sufficiently large universal constant $b$.
The particular circuit for a function need not be unique.

\begin{proposition}[Established analytic ceiling]\label{prop:ceiling}
There are universal positive constants $b,c_0,c_1$ such that the following
holds for every integer sample budget $n\ge1$ for the original posterior \eqref{eq:prior}--\eqref{eq:conditioning}.
Put $R=d+n+L+2$ and suppose
\begin{equation}\label{eq:padding}
 L\ge c_0\log R.
\end{equation}
For every $0<\beta<1$, every actual integer width
\begin{equation}\label{eq:width}
 N\ge 2^{c_1d}R^{c_1}+c_1(L-1)\log(1/\beta),
\end{equation}
every $0<t\le1$, and every consistent sample of at most $n$ observations
whose conditioning event has positive probability,
\begin{equation}\label{eq:ceiling}
 \Pi_S\{W:\sgn z_W\in\mathcal H_R\}\ge1-\beta.
\end{equation}
The class $\mathcal H_R$ is common to all these samples, widths, and
scales. In particular, the width threshold is independent of $t$ and
of the input law.
\end{proposition}

This is the circuit conclusion of the uniform actual-posterior ceiling
in \cite{barriers}, where it is labeled
\texttt{thm:uniform-ceiling}; no new analytic ceiling is claimed here. Its proof has three ingredients. First,
finite-event relative-entropy certificates associate most posterior mass
with controlled covariance paths of bounded action. A target-mass
construction for a sample memorizer supplies the posterior denominator
needed for this step. Second, the bounded-action paths admit finite
canonical feature representations, with bounded-size covers even when
$t$ approaches one. Third, parallel iteration of the resulting feature
recursion gives polynomial circuit size and logarithmic threshold depth.
The actual finite-width estimate, rather than only a limiting variational
statement, is essential to \eqref{eq:ceiling}.

The present proof uses Proposition~\ref{prop:ceiling} only with
$\beta=1/6$. It does not use the explicit covariance-path rate, an
optimal representation, or an efficient way of constructing a circuit
from network weights. Nor does it assert that one circuit carries most
posterior mass. The posterior can be an arbitrary mixture over its
ceiling class.

\section{From distribution-free mixtures to an exact circuit}
\label{sec:majorities}

The next lemma is independent of neural networks. It is an elementary
finite boosting argument. Its deterministic weight update avoids any
need for a minimax theorem or a concentration inequality.

\begin{lemma}[A small exact majority from weak prediction]
\label{lem:majority}
Let $X$ be a finite nonempty set with $|X|\le2^d$, let
$f:X\to\{-1,1\}$, and let $\mathcal H$ be a nonempty finite class of
Boolean functions. Suppose that for every probability law $\mu$ on $X$
there is an $h\in\mathcal H$ with
\[
 \mu\{x:h(x)\ne f(x)\}\le9/20.
\]
Then $f$ is exactly a majority of $O(d+1)$ members of $\mathcal H$,
allowing repetitions. In fact $400(d+1)$ voters suffice.
\end{lemma}
\begin{proof}
Initially give every input weight $w_0(x)=1$. At step $j$, normalize the
positive weights to a probability law and choose a voter $h_j$ whose
weighted error is at most $9/20$. Multiply a correctly classified input's
weight by $9/10$, and an incorrectly classified input's weight by $10/9$.
If $W_j=\sum_x w_j(x)$ is the total weight, then
\[
 \frac{W_{j+1}}{W_j}
 \le\frac{11}{20}\frac9{10}+\frac9{20}\frac{10}{9}
 =\frac{199}{200}.
\]
After $m$ rounds,
\begin{equation}\label{eq:boost-potential}
 W_m\le |X|(199/200)^m\le2^d e^{-m/200}.
\end{equation}
Set $m=400(d+1)$. For an elementary rational check, Bernoulli's inequality
gives $(200/199)^{400}\ge1+400/199>2$, so
$(199/200)^{400}<1/2$. Consequently $W_m\le1/2<1$.
But an input misclassified by at least half the voters would have final
weight
\[
 (9/10)^{\#\mathrm{correct}}(10/9)^{\#\mathrm{wrong}}\ge1,
\]
contradicting $W_m<1$. Every input is therefore correctly classified
by a strict majority of the chosen voters.
\end{proof}

\begin{lemma}[A successful posterior supplies a weak voter]
\label{lem:posterior-voter}
Let $\pi$ be a distribution on Boolean functions on $X$, and suppose
$\pi(\mathcal H)\ge5/6$. If
\[
 \mu\{x:\pi\{h:h(x)=f(x)\}\le2/3\}\le1/16,
\]
then some $h\in\mathcal H$ has $\mu$-error at most $9/20$.
\end{lemma}
\begin{proof}
At every point outside the confidence-failure set, a draw from $\pi$
has wrong-label probability less than $1/3$. On that set its error
probability is at most one. Averaging over inputs gives
\[
 \E_{h\sim\pi}\mu\{x:h(x)\ne f(x)\}
 \le\frac13+\frac23\frac1{16}=\frac38.
\]
Conditioning on $\mathcal H$ increases this average by at most $6/5$,
so its conditional average is at most $(3/8)(6/5)=9/20$. A finite
average cannot be smaller than all its contributing terms. One
$h\in\mathcal H$ therefore has error at most $9/20$.
\end{proof}

\begin{proposition}[Posterior success forces a small parallel circuit]
\label{prop:success-circuit}
Fix $d,n,L$ satisfying \eqref{eq:padding}, one width satisfying
\eqref{eq:width} with $\beta=1/6$, and one scale $0<t\le1$.
These choices are made before the input distribution.
Suppose that for every input law $\mu$, there is positive probability of
an iid labeled $n$-sample with $\mathcal E_\mu(f,S)\le1/16$.
Then $f$ has a threshold circuit of size polynomial in $d+n+L$ and depth
$O(\log(d+n+L+2))$. The constants are universal and independent of
learner width and scale.
\end{proposition}
\begin{proof}
For each law $\mu$, choose a successful sample and push its posterior
forward to Boolean functions. Proposition~\ref{prop:ceiling} gives mass
at least $5/6$ to the common finite class $\mathcal H_R$.
Lemma~\ref{lem:posterior-voter} supplies one member of that class with
$\mu$-error at most $9/20$.

This is precisely the distribution-free hypothesis of
Lemma~\ref{lem:majority}. Its exact majority has $O(d+1)$ voters from
$\mathcal H_R$. Compile the voters in parallel and append one majority
gate. Size grows by only a polynomial factor, and depth increases by one.
\end{proof}

The posterior mixture is handled by averaging; it is not identified with
one selected predictor. The successful samples and chosen voters may
vary with the input law, while $\mathcal H_R$ remains fixed. Positive
sample-success probability is used only to supply one successful sample
for each distribution. An earlier two-majority/minimax proof in
\cite{barriers} permits every confidence-failure tolerance below $1/2$;
the smaller fixed tolerance here makes the proof substantially shorter.

\section{Small Boolean circuits have bounded-weight teachers}

\begin{lemma}[Digital tanh implementation]\label{lem:teacher}
Every NAND circuit with $s$ gates on $d$ inputs has a tanh implementation
of hidden width $C=d+s+256$ and weighted depth $L=d+s+3$, with all
weights integers of absolute value at most $192$ and signed output
greater than $1/2$ throughout $X_d$. It has no hidden biases. The depth
can be further increased by inserting copy layers with the same bounds.
Other fixed bounded-fan-in Boolean bases incur only a constant-factor
change in circuit size before this construction.
\end{lemma}
\begin{proof}
Use signed Boolean wire values. NAND, NOT, and COPY have the scores
\[
 1-u-v,\qquad -u,\qquad u.
\]
Their signed margins are at least one on exact Boolean inputs.
Reserve coordinates for the $d$ inputs, the two Boolean constants,
and every computed gate. Initialize their tanh encodings in one hidden
layer from the augmented input. Then process one NAND gate per layer,
retaining all earlier values by copying. The true constant wire supplies
all subsequent affine constants. Append $d+1$ copy layers. The number of
hidden-to-hidden stages is $s+d+1$, hence the weighted depth is $s+d+3$.
The width $C=d+s+256$ leaves all required coordinates available.

Use gain $64$ at each gate. Suppose that each incoming wire is within
$1/16$ of its intended sign. The score perturbation, including the
constant wire, is at most $3/16$. Thus the signed preactivation is at
least $64/2$. The elementary saturation bound
\[
 by\ge A/2,\quad b\in\{-1,1\},\quad A>0
 \quad\Longrightarrow\quad |\tanh y-b|\le4/A
\]
then gives a new error at most $1/16$. For example, the bound follows
from $1-\tanh v=2/(e^{2v}+1)$ and $e^a\ge1+a$.
The initial input and constant signs satisfy the same estimate.
This proves an error invariant uniform in the number of layers.

A coefficient receives at most three contributions of magnitude $64$,
including coincident input and constant wires, so its absolute value is
at most $192$. The final linear readout is twice the selected output
wire. Its signed margin is at least $2(1-1/16)=15/8>1$.
Further copy layers preserve the invariant. Unused neurons have zero
weights, so zero padding to a larger width also preserves the function.
\end{proof}

The teacher construction is digitally stable at its Boolean gate
margins. The conditional negative result nevertheless concerns the
\emph{unrestricted learner prior}; having a stable implementation of the
target does not force the posterior's competing explanations to share
that stability.

\section{Proof of the main theorem}

We prove the contrapositive in its slightly stronger positive form.

\begin{theorem}[A universal polynomial guarantee would collapse classes]
\label{thm:collapse}
For every schedule \eqref{eq:schedule},
\[
 \mathsf{PolyLearn}(t)\quad\Longrightarrow\quad\Pp=\TC^1.
\]
\end{theorem}
\begin{proof}
Assume Definition~\ref{def:poly}, with fixed $A,k$, and choose a positive
integer $A_0\ge A$. Take any family
$(f_d)$ in $\Pp$, represented by Boolean circuits of size polynomial in
$d$. Lemma~\ref{lem:teacher} gives teachers with width, weight bound,
and depth polynomial in $d$. Their width $C_d$ is at least $256$, so it
also bounds every coefficient, and their depth $L_d$ is already at least
$d$.

The prescribed sample count
\[
 n_d=A_0(1+d+C_d+L_d)^k
\]
is polynomial in $d$ and is fixed \emph{before} choosing learner width.
Thus $R_d=d+n_d+L_d+2$ is polynomial in $d$. Since $L_d\ge d$, the
padding condition $L_d\ge c_0\log R_d$ holds for all sufficiently large
$d$.

For each such dimension choose one integer width at least $C_d$ and at
least the threshold \eqref{eq:width} with $\beta=1/6$. The threshold is
independent of $t$, so now set $t_d=t(d,L_d,N_d)$. There is no circular
choice of width and scale. Both are fixed before the input law.

The assumed learning guarantee gives success probability at least $1/2$
for every input law. Proposition~\ref{prop:success-circuit}, applied
with confidence-failure tolerance $1/16$, produces a threshold circuit for $f_d$ of
polynomial size in $R_d$, hence polynomial size in $d$, and depth
$O(\log R_d)=O(\log(d+2))$. The constants are uniform in $d$.
Finitely many remaining input dimensions can be hardcoded.

Every family in $\Pp$ therefore belongs to nonuniform $\TC^1$.
The reverse inclusion is standard, so $\Pp=\TC^1$.
\end{proof}

Taking the contrapositive proves Theorem~\ref{thm:main}.
For fixed $A,k$, a violation of \eqref{eq:guarantee} has probability
strictly less than $1/2$ for the event $\mathcal E_\mu(f,S)\le1/16$.
Its complementary event is exactly \eqref{eq:failure}. All distributions
and samples here are on a finite input set, so no measurability issue is
hidden in this logical step.

The reduction has the following order of choices:
\[
 \begin{aligned}
 &\text{sample polynomial}\ \longrightarrow\
 \text{target family and teachers}\ \longrightarrow\ n_d,\\
 &n_d\ \longrightarrow\ N_d\ \longrightarrow\ t_d
 \ \longrightarrow\ \text{input law and sample}.
 \end{aligned}
\]
The last two objects may vary while the architecture and scale remain
fixed. A sample bound polynomial in $N$ would not permit this argument:
it would make $n_d$ depend on the large width chosen only afterwards.

\subsection{A fixed hard family and one law for all sufficiently large widths}

The argument gives more than the failure of a particular confidence
threshold. The following consequence separates what can be made uniform
from the absence of an explicit asymptotic lower-bound rate.

\begin{corollary}[Failure for every sample on a hard distribution]
\label{cor:hard-family}
Assume $\Pp\ne\TC^1$. There is a fixed family $(f_d)$ in $\Pp$ with
bounded-weight tanh teachers of width $C_d=\poly(d)$ and depth
$d\le L_d=\poly(d)$ such that the following holds.
For every polynomially bounded integer sample budget $n_d\ge1$, there
are infinitely many dimensions $d$ with an input law $\mu_d$ for which
\[
 \mathcal E_{\mu_d}(f_d,S)>1/16
\]
for every labeled sample $S$ of length at most $n_d$, at every width above
both $C_d$ and the analytic ceiling threshold, and every scale $0<t\le1$.
The law $\mu_d$ is common to all those samples, widths, and scales.
In particular, the probability of failure is one for iid sampling from
$\mu_d$ at any of those sample counts.
\end{corollary}
\begin{proof}
Fix any family in $\Pp\setminus\TC^1$ and compile its teachers as in the
proof of Theorem~\ref{thm:collapse}. Put $R_d=d+n_d+L_d+2$.
For all sufficiently large $d$, the padding condition holds. Let
$\mathcal K_d$ be the class of all majorities of $400(d+1)$ members
of $\mathcal H_{R_d}$, allowing repetitions.
Every member has size polynomial in $R_d$ and depth $O(\log R_d)$.
Since $R_d=\poly(d)$, if $f_d\in\mathcal K_d$ for all but finitely many
$d$, the entire family would be in $\TC^1$. Hence $f_d\notin\mathcal K_d$
for infinitely many $d$.

At each such dimension, the contrapositive of
Lemma~\ref{lem:majority} gives one law $\mu_d$ for which every member of
$\mathcal H_{R_d}$ has error strictly greater than $9/20$.
Proposition~\ref{prop:ceiling} applies to every eligible width, scale,
and consistent sample of at most $n_d$ points, using the same class
$\mathcal H_{R_d}$. By Lemma~\ref{lem:posterior-voter}, none of those
posteriors can have confidence-failure mass at most $1/16$ under $\mu_d$.
This proves the claim. All target-labeled samples have positive
conditioning probability because the teacher exists.
\end{proof}

The teacher family is fixed before the sample budget. The hard input
law may depend on the budget and dimension. No particular growth rate
or explicit description of the hard family follows from the separation
assumption alone.

\section{Interpretation and limits}

\paragraph{What ``no polynomial bound'' means.}
The conclusion defeats every universal polynomial resource bound. It
does not assert an exponential lower bound in $L$, or even a fixed
stretched-exponential rate. Quasipolynomial or other intermediate sample
growth remains compatible with the theorem. The input dimension and
teacher resources vary together; fixing the finite input domain would
give an elementary finite-domain upper bound regardless of depth.

\paragraph{Why the input distribution is allowed to vary.}
For each proposed polynomial, the hard instance may use its own common
training/test law. The boosting argument would not follow from success
only under the uniform distribution. The theorem therefore supplies no
uniform-input lower bound without additional work.

\paragraph{No contradiction with positive fixed-depth theorems.}
For each fixed $L$, a bound with degree depending on $L$ can hold.
The theorem prohibits one degree working jointly over all depths and
target resources. It also leaves intact width-dependent finite-width
bounds and positive bounds for particular target representations.

\paragraph{Robust prior restrictions.}
Positive joint-polynomial results obtained by discarding unstable
networks change the learner. Theorem~\ref{thm:main} gives a conditional
reason that some change is needed for a theorem with the full original
target and distribution quantifiers. It does not identify the best
restriction, establish that trained networks satisfy one, or prove that
support restriction is preferable to every other modification.

\paragraph{Computational and statistical claims.}
The use of a circuit-complexity conjecture does not make this a lower
bound on the runtime of inference. The posterior is the exact conditional
Gaussian law. Circuit complexity enters because of the functions to
which that posterior assigns mass, and because distribution-free
prediction can be converted into an exact finite circuit.

\paragraph{Formal verification scope.}
The accompanying Lean project defines the actual Gaussian prior,
tanh network, global-margin posterior, iid sample law, and concrete
Boolean and threshold circuits used in the statement. It proves the
bounded-weight teacher compiler, the posterior-to-weak-voter transfer,
the deterministic boosting argument, the actual circuit composition,
threshold-to-Boolean simulation, polynomial family resource bounds,
and the final circuit-class implication.

There is one remaining analytic premise. The Lean theorem
\begin{center}
\texttt{conditionalImpossibility\_of\_posteriorCeiling}
\end{center}
proves the implication from the concrete proposition
\texttt{AnalyticPosteriorCeiling} to the concrete main claim
\texttt{ConditionalImpossibility}. The former is the finite-width
Gaussian posterior statement corresponding to
Proposition~\ref{prop:ceiling}, with integer resource envelopes; it is
not an abstract oracle standing in for a different learner.
\emph{Its proof has not been formalized.} The latter proposition states
$\Pp\ne\TC^1\Rightarrow\forall t\,\neg\mathsf{PolyLearn}(t)$ with
the exact conventions above, but has no closed Lean proof from the
complexity assumption alone. Neither proposition is postulated as an
axiom. Thus the package formally verifies the reduction and its
concrete definitions, not the complete main theorem. The verification
record gives file-level coverage and exact build information.

\begin{thebibliography}{9}
\bibitem{barriers}
ChatGPT-6 Astra, Kaarel H\"anni, and Dmitry Vaintrob.
\emph{Frobenius product, parallel circuits, and sample-complexity barriers
for deep Gaussian Bayesian networks}. Research manuscript, 2026.
\url{https://github.com/kaarelh/AI-works/tree/main/neural-network-generalization/works-in-progress/boolean-depth-barriers}.
The present paper uses the uniform actual-posterior circuit ceiling,
Theorem \texttt{thm:uniform-ceiling}, in the accompanying LaTeX source.
\end{thebibliography}
\end{document}
